\chapter{Capitolo 4 - Oggetti, morfismi, composizione, identità in estensione OCT}

\section{1. Problema}

Con i Capitoli 1-3 abbiamo costruito:

\begin{enumerate}
\def\labelenumi{\arabic{enumi}.}
\tightlist
\item
  la motivazione epistemica dell'estensione ordinativa;
\item
  la base ontologica della singolarità;
\item
  il kernel assiomatico O1-O7.
\end{enumerate}

Con i Capitoli 8-11 abbiamo consolidato:

\begin{enumerate}
\def\labelenumi{\arabic{enumi}.}
\tightlist
\item
  lo spazio di validità \texttt{S\_Omega};
\item
  il blocco teorematico fondativo e differenziale;
\item
  la grammatica di controesempi, limiti e falsificabilità.
\end{enumerate}

Ora inizia la fase \texttt{v1.3}: espansione sistematica del corpus categoriale classico in quadro OCT. Il primo passo è apparentemente semplice ma in realtà decisivo: ridefinire \textbf{oggetti, morfismi, composizione e identità} in forma compatibile con il classico ma dotata di selettività ordinativa.

Il problema specifico del capitolo è:

\textbf{come preservare la potenza della sintassi categoriale classica evitando, allo stesso tempo, che entità formalmente lecite ma ordinativamente sterili o degenerative vengano trattate come equivalenti?}

La difficoltà emerge su due fronti:

\begin{enumerate}
\def\labelenumi{\arabic{enumi}.}
\tightlist
\item
  se l'estensione è troppo debole, non aggiunge discriminazione reale;
\item
  se l'estensione è troppo invasiva, rompe la continuità con la teoria classica.
\end{enumerate}

La strategia adottata qui è conservativa-selettiva:

\begin{enumerate}
\def\labelenumi{\arabic{enumi}.}
\tightlist
\item
  conservativa sul piano strutturale;
\item
  selettiva sul piano ordinativo.
\end{enumerate}

In altre parole:

\begin{enumerate}
\def\labelenumi{\arabic{enumi}.}
\tightlist
\item
  non cambiamo il significato classico di composizione e identità;
\item
  introduciamo un livello di ammissibilità che decide quando la stessa struttura è operativamente reale in \texttt{Omega}.
\end{enumerate}

\subsection{1.1 Perché questo capitolo viene dopo la fase 8-11}

Ordine non lineare ma intenzionale:

\begin{enumerate}
\def\labelenumi{\arabic{enumi}.}
\tightlist
\item
  prima abbiamo fissato lo spazio di validità e i teoremi;
\item
  ora torniamo alla micro-struttura classica con criteri già stabilizzati.
\end{enumerate}

Il vantaggio è metodologico:

\begin{enumerate}
\def\labelenumi{\arabic{enumi}.}
\tightlist
\item
  le definizioni qui introdotte non sono ``promesse'', ma già agganciate a metriche e criteri di fallimento;
\item
  ogni scelta locale su oggetto/morfismo/composizione/identità può essere testata contro il quadro critico del Capitolo 11.
\end{enumerate}

\subsection{1.2 Obiettivo tecnico del capitolo}

L'obiettivo non è riscrivere da zero la teoria degli oggetti e dei morfismi, ma esplicitare:

\begin{enumerate}
\def\labelenumi{\arabic{enumi}.}
\tightlist
\item
  qual è il corrispettivo OCT di ogni nozione classica;
\item
  quali condizioni aggiuntive servono per ammissibilità ordinativa;
\item
  quali errori di inferenza sono evitati dalla riformulazione.
\end{enumerate}

\begin{center}\rule{0.5\linewidth}{0.5pt}\end{center}

\section{2. Tesi del capitolo}

La tesi del capitolo è articolata in otto enunciati.

\begin{enumerate}
\def\labelenumi{\arabic{enumi}.}
\tightlist
\item
  L'oggetto ordinativo è una singolarizzazione contestuale di un oggetto classico, non una sostituzione antagonista.
\item
  Il morfismo ordinativo è un morfismo classico ben tipizzato più vincolo di ammissibilità \texttt{Coh/Phi}.
\item
  La composizione ordinativa è chiusa solo condizionatamente (gate sulla composta).
\item
  L'identità ordinativa richiede neutralità non-degenerativa, non sola tautologia sintattica.
\item
  La struttura classica è preservata dal funtore di oblio, ma non vale il viceversa in generale.
\item
  Le quattro nozioni classiche restano formalmente intatte e vengono rafforzate da criteri operativi.
\item
  La distinzione tra validità locale e validità composita è obbligatoria in OCT.
\item
  Il capitolo prepara direttamente limiti/colimiti/equivalenze del Capitolo 5.
\end{enumerate}

\begin{center}\rule{0.5\linewidth}{0.5pt}\end{center}

\section{3. Definizioni canoniche}

\subsection{Definizione 4.1 (Oggetto ordinativo)}

Dato un oggetto classico \texttt{A}, definiamo l'oggetto ordinativo in \texttt{Omega}:

\texttt{A\_Omega\ :=\ S\_Omega(A)\ =\ \textless{}A,\ R\_Omega(A),\ H\_t(A),\ Pi\_Omega(A)\textgreater{}}.

Commento:

\begin{enumerate}
\def\labelenumi{\arabic{enumi}.}
\tightlist
\item
  \texttt{A} resta il supporto strutturale;
\item
  \texttt{R\_Omega}, \texttt{H\_t}, \texttt{Pi\_Omega} introducono la specificazione relazionale-funzionale minima.
\end{enumerate}

\subsection{Definizione 4.2 (Morfismo ordinativamente ammissibile)}

Per \texttt{f:\ A\_Omega\ -\textgreater{}\ B\_Omega}:

\texttt{Adm\_Omega(f)\ :=\ Syn(f)\ and\ Coh\_Omega(f)\ \textgreater{}=\ tau\_f\ and\ Phi\_Omega(f)\ !=\ 0}.

Commento:

\begin{enumerate}
\def\labelenumi{\arabic{enumi}.}
\tightlist
\item
  \texttt{Syn(f)} è requisito necessario;
\item
  \texttt{Coh/Phi} sono requisiti selettivi.
\end{enumerate}

\subsection{Definizione 4.3 (Composizione ordinativa)}

Per \texttt{f:\ A\_Omega-\textgreater{}B\_Omega}, \texttt{g:\ B\_Omega-\textgreater{}C\_Omega}:

\texttt{g\ \ o\ \_Omega\ f} è ammessa se:

\begin{enumerate}
\def\labelenumi{\arabic{enumi}.}
\tightlist
\item
  \texttt{g\ o\ f} è definita in senso classico;
\item
  \texttt{Adm\_Omega(f)} e \texttt{Adm\_Omega(g)};
\item
  \texttt{Adm\_Omega(g\ o\ f)}.
\end{enumerate}

\subsection{Definizione 4.4 (Identità ordinativa)}

\texttt{id\_A\^{}Omega:\ A\_Omega\ -\textgreater{}\ A\_Omega} è identità ordinativa se:

\begin{enumerate}
\def\labelenumi{\arabic{enumi}.}
\tightlist
\item
  coincide con identità classica sul supporto \texttt{A};
\item
  non induce degradazione oltre soglia su coerenza/funzione in auto-applicazione.
\end{enumerate}

\subsection{Definizione 4.5 (Classe di composizione sterile)}

Una composizione è sterile se:

\begin{enumerate}
\def\labelenumi{\arabic{enumi}.}
\tightlist
\item
  è sintatticamente lecita;
\item
  mantiene coerenza sopra soglia;
\item
  ma \texttt{Phi\_Omega=0}.
\end{enumerate}

Questa classe è distinta dalla degenerazione pura (\texttt{Coh\textless{}tau}).

\subsection{Definizione 4.6 (Funtore di oblio ordinativo)}

\texttt{U\_Omega:\ OrdCat\_Omega\ -\textgreater{}\ Cat\_class}

rimuove il layer ordinativo mantenendo oggetti, morfismi e leggi classiche.

\begin{center}\rule{0.5\linewidth}{0.5pt}\end{center}

\section{4. Sviluppo formale}

\subsection{4.1 Oggetti: dal nodo astratto al nodo tipizzato}

Nel quadro classico, un oggetto è identificato dal suo ruolo in una rete di morfismi. Questa impostazione resta valida in OCT, ma viene resa operativa con un vincolo contestuale.

Passaggio:

\begin{enumerate}
\def\labelenumi{\arabic{enumi}.}
\tightlist
\item
  classico: \texttt{A};
\item
  ordinativo: \texttt{A\_Omega\ =\ \textless{}A,R,H,Pi\textgreater{}}.
\end{enumerate}

Questa estensione non altera le leggi classiche. Introduce però una differenza cruciale:

\begin{enumerate}
\def\labelenumi{\arabic{enumi}.}
\tightlist
\item
  due oggetti possono essere formalmente equivalenti in un livello astratto;
\item
  ma divergere ordinativamente per profilo relazionale, storia e funzione emergente.
\end{enumerate}

Errore evitato:

\begin{enumerate}
\def\labelenumi{\arabic{enumi}.}
\tightlist
\item
  assumere che identità simbolica o equivalenza strutturale basti per identità operativa.
\end{enumerate}

\subsection{4.2 Morfismi: tipizzazione più generatività}

In category theory classica, la validità del morfismo è primariamente tipologica. In OCT questo resta il primo filtro, ma non l'ultimo.

Schema:

\begin{enumerate}
\def\labelenumi{\arabic{enumi}.}
\tightlist
\item
  filtro 1: \texttt{Syn(f)} (classico);
\item
  filtro 2: \texttt{Coh\_Omega(f)\textgreater{}=tau\_f};
\item
  filtro 3: \texttt{Phi\_Omega(f)!=0}.
\end{enumerate}

Ne segue una distinzione esplicita:

\begin{enumerate}
\def\labelenumi{\arabic{enumi}.}
\tightlist
\item
  morfismi leciti ma sterili;
\item
  morfismi leciti ma degenerativi;
\item
  morfismi ordinativamente ammissibili.
\end{enumerate}

Questa tripartizione migliora la diagnosi locale e prepara la composizione controllata.

\subsection{4.3 Composizione: chiusura condizionata}

La legge classica \texttt{g\ o\ f} resta intatta. Quello che cambia è la nozione di ``chiusura utile''.

In OCT:

\begin{enumerate}
\def\labelenumi{\arabic{enumi}.}
\tightlist
\item
  non basta che \texttt{f} e \texttt{g} siano ammissibili singolarmente;
\item
  serve test ordinativo sulla composta.
\end{enumerate}

Questo evita inferenze lineari errate:

\begin{enumerate}
\def\labelenumi{\arabic{enumi}.}
\tightlist
\item
  ``due passi validi implicano sempre traiettoria valida''.
\end{enumerate}

Controesempio tipico:

\begin{enumerate}
\def\labelenumi{\arabic{enumi}.}
\tightlist
\item
  \texttt{f} preserva coerenza locale;
\item
  \texttt{g} preserva coerenza locale;
\item
  \texttt{g\ o\ f} amplifica incompatibilità latenti e cade sotto soglia.
\end{enumerate}

Conclusione:

\begin{enumerate}
\def\labelenumi{\arabic{enumi}.}
\tightlist
\item
  la composizione ordinativa è una legge con gate, non solo una concatenazione.
\end{enumerate}

\subsection{4.4 Identità: neutralità condizionata}

Nel classico, \texttt{id\_A} è neutra per definizione strutturale. In OCT la neutralità è anche una proprietà misurabile.

Richiesta minima:

\begin{enumerate}
\def\labelenumi{\arabic{enumi}.}
\tightlist
\item
  auto-applicazione non produce deriva sistematica su \texttt{Coh/Phi}.
\end{enumerate}

Perché è necessario:

\begin{enumerate}
\def\labelenumi{\arabic{enumi}.}
\tightlist
\item
  in sistemi dinamici reali, un'operazione formalmente neutra può attivare costi nascosti, attriti di sincronizzazione o perdita funzionale.
\end{enumerate}

Errore evitato:

\begin{enumerate}
\def\labelenumi{\arabic{enumi}.}
\tightlist
\item
  confondere neutralità simbolica con neutralità operativa.
\end{enumerate}

\subsection{4.5 Composizione locale vs composizione globale}

Distinzione obbligatoria:

\begin{enumerate}
\def\labelenumi{\arabic{enumi}.}
\tightlist
\item
  validità locale: test su singolo morfismo;
\item
  validità globale: test su catena compositiva.
\end{enumerate}

In OCT entrambe sono necessarie:

\begin{enumerate}
\def\labelenumi{\arabic{enumi}.}
\tightlist
\item
  la sola validità locale non garantisce tenuta globale;
\item
  la sola validità globale senza tracciamento locale rende opaca la diagnosi di fallimento.
\end{enumerate}

Questa distinzione è fondamentale per pipeline modulari e sistemi multi-stadio.

\subsection{4.6 Oggetti e morfismi come coppia inseparabile}

Nel quadro ordinativo, non ha senso valutare un oggetto senza dinamica morfismica, né un morfismo senza contesto oggettuale singolarizzato.

Principio operativo:

\begin{enumerate}
\def\labelenumi{\arabic{enumi}.}
\tightlist
\item
  \texttt{A\_Omega} determina condizioni di ammissibilità dei morfismi uscenti/entranti;
\item
  l'insieme dei morfismi ammissibili retro-definisce lo stato operativo di \texttt{A\_Omega}.
\end{enumerate}

Ne deriva una co-definizione dinamica:

\begin{enumerate}
\def\labelenumi{\arabic{enumi}.}
\tightlist
\item
  oggetti e morfismi sono analiticamente distinguibili ma operativamente co-determinati.
\end{enumerate}

\subsection{4.7 Stabilità di composizione in vicinato contestuale}

Per prevenire fragilità da fitting locale, la composizione va valutata anche su \texttt{N\_epsilon(Omega)}.

Richiesta minima:

\begin{enumerate}
\def\labelenumi{\arabic{enumi}.}
\tightlist
\item
  \texttt{Adm\_Omega(g\ o\ f)} deve restare stabile entro perturbazioni ammesse.
\end{enumerate}

In assenza di questa condizione:

\begin{enumerate}
\def\labelenumi{\arabic{enumi}.}
\tightlist
\item
  si ottengono composizioni ``formalmente vincenti'' ma non replicabili.
\end{enumerate}

\subsection{4.8 Relazione con la matrice A/B/C/D}

Oggetti e morfismi in OCT possono essere classificati tramite la matrice standard:

\begin{enumerate}
\def\labelenumi{\arabic{enumi}.}
\tightlist
\item
  A: ammissibile pieno;
\item
  B: lecito ma sterile;
\item
  C: lecito ma degenerativo;
\item
  D: non lecito sintatticamente.
\end{enumerate}

Nel Capitolo 4 questa matrice è utile per:

\begin{enumerate}
\def\labelenumi{\arabic{enumi}.}
\tightlist
\item
  qualificare rapidamente lo stato di un morfismo;
\item
  spiegare perché una composizione viene accettata o rifiutata.
\end{enumerate}

\subsection{4.9 Recupero classico come controllo di coerenza}

Il capitolo conserva esplicitamente la riduzione classica:

\begin{enumerate}
\def\labelenumi{\arabic{enumi}.}
\tightlist
\item
  se il layer ordinativo è disattivato, restano oggetti/morfismi/composizione/identità classici.
\end{enumerate}

Questo punto impedisce due letture sbagliate:

\begin{enumerate}
\def\labelenumi{\arabic{enumi}.}
\tightlist
\item
  OCT come negazione del classico;
\item
  OCT come lessico alternativo senza nuove regole.
\end{enumerate}

\subsection{4.10 Mappa errori tipici evitati dal capitolo}

Errori classici in applicazione ingenua:

\begin{enumerate}
\def\labelenumi{\arabic{enumi}.}
\tightlist
\item
  ``tipizzato quindi valido'';
\item
  ``composto quindi robusto'';
\item
  ``identità quindi neutro in ogni senso'';
\item
  ``equivalenza locale quindi equivalenza globale''.
\end{enumerate}

Il capitolo fornisce, per ciascun errore, un criterio di controllo operativo.

\subsection{4.11 Composizione lunga e accumulo di rischio ordinativo}

Una catena compositiva lunga \texttt{f\_n\ \ o\ \ ...\ \ o\ \ f\_1} introduce un problema che il controllo locale non risolve da solo: l'accumulo di rischio.

Anche quando ogni \texttt{f\_i} soddisfa il gate locale, il sistema può mostrare:

\begin{enumerate}
\def\labelenumi{\arabic{enumi}.}
\tightlist
\item
  drift progressivo di coerenza;
\item
  compressione della funzione emergente;
\item
  amplificazione di perdite marginali in punti di interfaccia.
\end{enumerate}

Per questo proponiamo un criterio di monitoraggio a due scale:

\begin{enumerate}
\def\labelenumi{\arabic{enumi}.}
\tightlist
\item
  \textbf{scala locale}: \texttt{Adm\_Omega(f\_i)} per ogni passo;
\item
  \textbf{scala cumulativa}: \texttt{Adm\_Omega(F\_k)} per ogni prefisso \texttt{F\_k\ =\ f\_k\ \ o\ \ ...\ \ o\ \ f\_1}.
\end{enumerate}

Il vantaggio è diagnostico:

\begin{enumerate}
\def\labelenumi{\arabic{enumi}.}
\tightlist
\item
  si individua il punto in cui la catena passa da robusta a fragile;
\item
  si evita di attribuire il collasso solo all'ultimo morfismo osservato.
\end{enumerate}

Questa analisi è particolarmente importante per pipeline IA multi-step e trasformazioni semantiche iterabili.

\subsection{4.12 Criterio di margine compositivo}

Oltre al gate binario (ammesso/non ammesso), introduciamo un margine compositivo:

\texttt{m\_comp(F\_k)\ :=\ Coh\_Omega(F\_k)\ -\ tau\_chain}.

Uso operativo:

\begin{enumerate}
\def\labelenumi{\arabic{enumi}.}
\tightlist
\item
  \texttt{m\_comp\ \textgreater{}\textgreater{}\ 0}: catena robusta;
\item
  \texttt{m\_comp\ \textgreater{}\ 0} ma piccolo: catena fragile, da monitorare;
\item
  \texttt{m\_comp\ \textless{}=\ 0}: catena degenerativa.
\end{enumerate}

Il margine permette decisioni più fini:

\begin{enumerate}
\def\labelenumi{\arabic{enumi}.}
\tightlist
\item
  non tutte le composizioni ammesse hanno la stessa affidabilità;
\item
  il valore di margine guida priorità di revisione e strategie di riprogettazione.
\end{enumerate}

In assenza di margine, si rischia una logica ``on/off'' che non cattura la dinamica di prossimità al collasso.

\subsection{4.13 Interfacce tra morfismi: punto critico nascosto}

Molti fallimenti globali non nascono dal singolo morfismo, ma dalla giunzione tra morfismi consecutivi. In termini pratici, la compatibilità tipologica non garantisce compatibilità ordinativa di interfaccia.

Definiamo quindi il concetto di \textbf{interfaccia ordinativa} tra \texttt{f} e \texttt{g}:

\begin{enumerate}
\def\labelenumi{\arabic{enumi}.}
\tightlist
\item
  allineamento semantico delle uscite/ingressi;
\item
  assenza di perdita sistematica di coerenza al passaggio;
\item
  mantenimento della funzione emergente in transizione.
\end{enumerate}

Un criterio minimo è:

\texttt{Int\_Omega(g,f)\ =\ true} se il passaggio \texttt{B\_f\ -\textgreater{}\ B\_g} non introduce delta critico su \texttt{Coh} o annullamento su \texttt{Phi}.

Conseguenze operative:

\begin{enumerate}
\def\labelenumi{\arabic{enumi}.}
\tightlist
\item
  debug mirato su confini di modulo;
\item
  migliore separazione tra errore ``di modulo'' ed errore ``di integrazione''.
\end{enumerate}

\subsection{4.14 Tipologia dei fallimenti su oggetti e morfismi}

Per rendere il capitolo utilizzabile nei cicli di validazione, proponiamo una tipologia minima dei fallimenti.

\begin{enumerate}
\def\labelenumi{\arabic{enumi}.}
\item
  \textbf{F-Obj1 (Oggetto sottospecificato)}\\
  \texttt{A\_Omega} non ha profilo relazionale sufficiente per test affidabile.
\item
  \textbf{F-Mor1 (Morfismo sterile)}\\
  \texttt{Syn} e \texttt{Coh} validi, ma \texttt{Phi=0}.
\item
  \textbf{F-Mor2 (Morfismo degenerativo)}\\
  \texttt{Syn} valido, ma \texttt{Coh\textless{}tau}.
\item
  \textbf{F-Comp1 (Composta fragile)}\\
  passi locali validi, composta al bordo di soglia.
\item
  \textbf{F-Comp2 (Composta collassata)}\\
  composta in classe degenerativa.
\item
  \textbf{F-Id1 (Identità non neutra)}\\
  auto-applicazione con drift ordinativo misurabile.
\end{enumerate}

Questa tipologia ha tre vantaggi:

\begin{enumerate}
\def\labelenumi{\arabic{enumi}.}
\tightlist
\item
  uniforma il lessico nei report di test;
\item
  collega direttamente tipo di fallimento e strategia di correzione;
\item
  riduce conflitti interpretativi tra team.
\end{enumerate}

\subsection{4.15 Criterio minimo di accettazione per modulo compositivo}

Proponiamo una regola di accettazione pratica per moduli OCT-based:

un modulo può essere dichiarato ``ordinativamente affidabile'' solo se soddisfa contemporaneamente:

\begin{enumerate}
\def\labelenumi{\arabic{enumi}.}
\tightlist
\item
  assenza di fallimenti F-Mor2 su suite primaria;
\item
  frequenza F-Mor1 sotto soglia dichiarata;
\item
  nessun caso F-Comp2 su scenari di carico nominale;
\item
  margine compositivo medio positivo sopra valore minimo dichiarato;
\item
  log completo delle interfacce critiche.
\end{enumerate}

Questa regola non pretende universalità assoluta, ma fornisce un baseline replicabile per sviluppo e pubblicazione tecnica.

\begin{center}\rule{0.5\linewidth}{0.5pt}\end{center}

\section{5. Proposizioni e teoremi locali}

\subsection{Proposizione 4.1 (Conservazione sintattica degli oggetti)}

Enunciato:
La singolarizzazione \texttt{S\_Omega} non distrugge la tipizzazione classica dell'oggetto.

Intuizione:
\texttt{A\_Omega} contiene \texttt{A} come supporto strutturale, con estensione relazionale-funzionale.

Stato: \texttt{validated}.

\subsection{Proposizione 4.2 (Necessità del gate morfismico)}

Enunciato:
Esistono morfismi \texttt{Syn(f)} veri ma \texttt{Adm\_Omega(f)} falsi.

Intuizione:
controesempi con \texttt{Phi=0} o \texttt{Coh\textless{}tau} rendono il gate indispensabile.

Stato: \texttt{validated}.

\subsection{Teorema 4.3 (Chiusura condizionata della composizione)}

Ipotesi:

\begin{enumerate}
\def\labelenumi{\arabic{enumi}.}
\tightlist
\item
  \texttt{Adm\_Omega(f)} e \texttt{Adm\_Omega(g)};
\item
  composizione classica definita;
\item
  \texttt{Adm\_Omega(g\ o\ f)} verificata.
\end{enumerate}

Tesi:
\texttt{g\ o\ \_Omega\ f} è composizione ordinativa valida.

Proof sketch:

\begin{enumerate}
\def\labelenumi{\arabic{enumi}.}
\tightlist
\item
  la definibilità classica garantisce coerenza sintattica;
\item
  il gate sulla composta garantisce non-degenerazione e non-sterilità;
\item
  dunque la composta è ammessa in \texttt{OrdCat\_Omega}.
\end{enumerate}

Conseguenze:
stabilisce il meccanismo operativo per catene compositive affidabili.

Stato: \texttt{in\_proof}.

\subsection{Teorema 4.4 (Neutralità ordinativa condizionata dell'identità)}

Ipotesi:

\begin{enumerate}
\def\labelenumi{\arabic{enumi}.}
\tightlist
\item
  \texttt{id\_A} classica su \texttt{A};
\item
  test ordinativo su auto-applicazione in \texttt{Omega};
\item
  assenza di degradazione oltre soglia.
\end{enumerate}

Tesi:
\texttt{id\_A\^{}Omega} è neutra in senso ordinativo.

Proof sketch:

\begin{enumerate}
\def\labelenumi{\arabic{enumi}.}
\tightlist
\item
  la neutralità sintattica è data;
\item
  la neutralità ordinativa richiede stabilità su \texttt{Coh/Phi};
\item
  soddisfatte entrambe, l'identità è pienamente neutra.
\end{enumerate}

Conseguenze:
evita uso improprio dell'identità in domini dinamici ad attrito.

Stato: \texttt{in\_proof}.

\subsection{Proposizione 4.5 (Non equivalenza locale-globale)}

Enunciato:
Validità ordinativa locale dei morfismi non implica validità ordinativa globale della catena.

Intuizione:
la composta può introdurre incompatibilità emergenti non visibili localmente.

Stato: \texttt{validated}.

\subsection{Teorema 4.6 (Accumulo non lineare del rischio compositivo)}

Ipotesi:

\begin{enumerate}
\def\labelenumi{\arabic{enumi}.}
\tightlist
\item
  catena compositiva \texttt{F\_n\ =\ f\_n\ \ o\ \ ...\ \ o\ \ f\_1} con \texttt{Adm\_Omega(f\_i)} veri;
\item
  esiste almeno un sottoinsieme di interfacce con \texttt{Int\_Omega} debole;
\item
  il contributo di rischio di interfaccia è non nullo e persistente.
\end{enumerate}

Tesi:
La validità locale uniforme non garantisce, in generale, margine compositivo globale stabile.

Proof sketch:

\begin{enumerate}
\def\labelenumi{\arabic{enumi}.}
\tightlist
\item
  ogni \texttt{f\_i} può rispettare il gate locale;
\item
  le interfacce deboli introducono perdite incrementali;
\item
  le perdite possono accumularsi oltre il bordo di sicurezza della catena;
\item
  ne segue che il margine globale può diventare fragile o negativo.
\end{enumerate}

Conseguenze:

\begin{enumerate}
\def\labelenumi{\arabic{enumi}.}
\tightlist
\item
  giustifica monitoraggio su prefissi compositivi;
\item
  fonda la necessità di metriche di interfaccia oltre le metriche di modulo.
\end{enumerate}

Stato: \texttt{in\_proof}.

\begin{center}\rule{0.5\linewidth}{0.5pt}\end{center}

\section{6. Implicazioni operative}

\subsection{6.1 Per implementazione di pipeline compositive}

Ogni pipeline OCT dovrebbe verificare:

\begin{enumerate}
\def\labelenumi{\arabic{enumi}.}
\tightlist
\item
  gate locale su ogni morfismo;
\item
  gate globale su composte critiche;
\item
  classificazione A/B/C/D per diagnosticare i fallimenti.
\end{enumerate}

\subsection{6.2 Per la documentazione di progetto}

È utile che ogni modulo espliciti:

\begin{enumerate}
\def\labelenumi{\arabic{enumi}.}
\tightlist
\item
  oggetti singolarizzati coinvolti;
\item
  morfismi ammessi e non ammessi;
\item
  ragione del rifiuto (sterilità vs degenerazione).
\end{enumerate}

Questo rende tracciabile il comportamento del sistema e accelera revisione.

\subsection{6.3 Per la validazione empirica}

I test dovrebbero includere:

\begin{enumerate}
\def\labelenumi{\arabic{enumi}.}
\tightlist
\item
  casi positivi di composizione robusta;
\item
  casi negativi di composizione fragile;
\item
  stress su identità in presenza di perturbazioni leggere.
\end{enumerate}

\subsection{6.4 Per il passaggio al Capitolo 5}

Limiti, colimiti, equivalenze e isomorfismi del Capitolo 5 useranno direttamente le nozioni qui fissate:

\begin{enumerate}
\def\labelenumi{\arabic{enumi}.}
\tightlist
\item
  oggetto ordinativo;
\item
  morfismo ammissibile;
\item
  composizione con gate;
\item
  identità non-degenerativa.
\end{enumerate}

\subsection{6.5 Per la governance editoriale}

Ogni nuova versione dovrebbe mantenere coerenza su questi quattro blocchi. Se cambia uno, devono essere riesaminati anche gli altri per evitare incoerenze interne nel manoscritto.

Inoltre, è consigliabile mantenere un changelog locale per oggetti, morfismi, composizioni e identità, così da collegare ogni modifica teorica a effetti misurabili sui benchmark.

\begin{center}\rule{0.5\linewidth}{0.5pt}\end{center}

\section{7. Limiti e non-claim}

Limiti espliciti:

\begin{enumerate}
\def\labelenumi{\arabic{enumi}.}
\tightlist
\item
  il capitolo non chiude ancora la formalizzazione completa dei casi enriched/2-categoriali;
\item
  la scelta delle soglie resta dominio-dipendente;
\item
  la neutralità ordinativa dell'identità richiede protocolli di misura robusti;
\item
  la distinzione locale-globale è formalizzata in versione minima, da estendere.
\end{enumerate}

Failure mode riconosciuti:

\begin{enumerate}
\def\labelenumi{\arabic{enumi}.}
\tightlist
\item
  sovra-ottimizzare i gate su dati specifici;
\item
  confondere casi sterili con casi degenerativi;
\item
  applicare il gate globale senza diagnosi locale;
\item
  dichiarare neutralità identitaria senza test ordinativo.
\end{enumerate}

Box C - Non-claim

Questo capitolo non implica:

\begin{enumerate}
\def\labelenumi{\arabic{enumi}.}
\tightlist
\item
  che ogni struttura classica debba essere automaticamente reinterpretata in OCT;
\item
  che il gate ordinativo elimini la necessità di prove matematiche classiche;
\item
  che l'ammissibilità ordinativa locale sia sufficiente per garantire robustezza sistemica.
\end{enumerate}

\begin{center}\rule{0.5\linewidth}{0.5pt}\end{center}

\section{8. Claim Ledger (S0-S3)}

{\def\LTcaptype{none} % do not increment counter
\begin{longtable}[]{@{}lllll@{}}
\toprule\noalign{}
ID & Claim & Tipo & Evidenza & Stato \\
\midrule\noalign{}
\endhead
\bottomrule\noalign{}
\endlastfoot
C4.1 & L'oggetto ordinativo conserva il supporto classico e aggiunge specificazione relazionale-funzionale & formale & S2 & in\_review \\
C4.2 & Il morfismo ordinativo richiede gate \texttt{Syn+Coh+Phi} & formale & S1 & validated \\
C4.3 & La composizione ordinativa è chiusa in modo condizionato e non automatico & formale & S2 & in\_proof \\
C4.4 & L'identità ordinativa richiede neutralità non-degenerativa verificata & formale & S2 & in\_proof \\
C4.5 & La validità locale non implica validità globale & formale-operativo & S1 & validated \\
C4.6 & La matrice A/B/C/D migliora la diagnosi operativa su oggetti e morfismi & metodologico & S1 & in\_review \\
\end{longtable}
}

\begin{center}\rule{0.5\linewidth}{0.5pt}\end{center}

\section{9. Bridge al Capitolo 5}

Il Capitolo 4 ha stabilito la grammatica micro-strutturale dell'estensione OCT:

\begin{enumerate}
\def\labelenumi{\arabic{enumi}.}
\tightlist
\item
  cosa sono gli oggetti in forma singolarizzata;
\item
  cosa rende un morfismo ammissibile;
\item
  quando una composizione è realmente valida;
\item
  cosa significa identità in senso non-degenerativo.
\end{enumerate}

Il Capitolo 5 userà questa base per affrontare il livello successivo:

\begin{enumerate}
\def\labelenumi{\arabic{enumi}.}
\tightlist
\item
  isomorfismi ed equivalenze;
\item
  limiti e colimiti;
\item
  selettività ordinativa delle costruzioni universali.
\end{enumerate}

In sintesi: il Capitolo 4 fissa i mattoni; il Capitolo 5 valuterà la stabilità delle architetture costruite con quei mattoni.

\begin{center}\rule{0.5\linewidth}{0.5pt}\end{center}

\section{Riferimenti}

Riferimenti di framework interno:

\begin{enumerate}
\def\labelenumi{\arabic{enumi}.}
\tightlist
\item
  Ghioni, F. (2026). \texttt{TE\_CORE\_v5.1.md}.
\item
  Ghioni, F. (2026). \texttt{Ordinative\_Set\_Theory\_OST\_A\_Concise\_Guide\_For\_AI\_v2\_1.md}.
\item
  \texttt{OCT\_FOUNDATIONAL\_BOOK\_CHAPTER\_03\_v1\_0.md}.
\item
  \texttt{OCT\_FOUNDATIONAL\_BOOK\_CHAPTER\_08\_v1\_0.md}.
\item
  \texttt{OCT\_FOUNDATIONAL\_BOOK\_CHAPTER\_09\_v1\_0.md}.
\item
  \texttt{OCT\_FOUNDATIONAL\_BOOK\_CHAPTER\_11\_v1\_0.md}.
\end{enumerate}

Riferimenti primari:

\begin{enumerate}
\def\labelenumi{\arabic{enumi}.}
\tightlist
\item
  Eilenberg, S., Mac Lane, S. (1945). \emph{General Theory of Natural Equivalences}.
\item
  Mac Lane, S. (1998). \emph{Categories for the Working Mathematician}.
\item
  Riehl, E. (2016). \emph{Category Theory in Context}.
\item
  Spivak, D. (2014). \emph{Category Theory for the Sciences}.
\end{enumerate}
